\chapter{Metodologia de desenvolvimento do projeto}
\label{cap:metodologia}

Este capítulo descreve como a equipe conduz o projeto: o processo seguido, a organização do trabalho no repositório, a gestão por indicadores PSM/CID, os cuidados éticos e o uso de ferramentas de inteligência artificial.

\section{Processo e organização do trabalho}
\label{sec:processo-trabalho}

A avaliação segue as quatro fases do processo de avaliação de produto (\autoref{fig:processo}), com uma entrega a cada duas fases \cite{ramos2026roteiro}. A equipe atua como \textbf{avaliador independente} (terceira parte): não desenvolve nem opera o SIGAA e não tem acesso a código-fonte ou infraestrutura, o que torna a avaliação externa (caixa-preta).

O trabalho é organizado no repositório Git da disciplina. Cada fase é desenvolvida em um ramo próprio (\texttt{fase-1}, \texttt{fase-2}) e integrada ao ramo principal por \textit{pull request} (PR), após revisão por outro membro. As mensagens de \textit{commit} recebem um identificador de tarefa (\texttt{[D01]} a \texttt{[D05]}), o que permite rastrear cada alteração até a atividade que a originou. A documentação é escrita em Markdown e publicada automaticamente como GitPage (MkDocs) a cada integração no ramo principal; cada página traz um histórico de versões com data, descrição e autor. Este artigo é mantido em LaTeX na pasta \texttt{artigo/} do mesmo repositório.

As decisões de cada fase são tomadas em reunião: cada membro elabora sua proposta individualmente, as propostas são consolidadas e as divergências são discutidas até o consenso. Esse procedimento foi usado, por exemplo, na atribuição das notas da matriz Impacto~$\times$~Risco (\autoref{sec:priorizacao}).

\section{Gestão do projeto com indicadores PSM/CID}
\label{sec:gestao}

O roteiro exige o acompanhamento do projeto com métricas do PSM/CID nas categorias calendário e progresso, recursos e custos e desempenho dos processos \cite{ramos2026roteiro}. A \autoref{tab:psm} apresenta os indicadores adotados, construídos segundo o modelo de informação do PSM \cite{mcgarry2002} e adaptados das especificações do \textit{framework} CID \cite{psmcid2}. Todos usam dados que o próprio repositório registra (\textit{commits}, PRs, \textit{issues} e marcos) ou uma planilha de esforço preenchida pelos membros.

\begin{table}[htb]
	\caption{Indicadores PSM/CID para a gestão do projeto}
	\label{tab:psm}
	\centering
	\scriptsize
	\begin{tabularx}{\textwidth}{L{1.9cm} L{3.2cm} L{3.6cm} Y}
		\toprule
		\textbf{Categoria} & \textbf{Necessidade de informação} & \textbf{Indicador (medidas)} & \textbf{Critério de decisão} \\
		\midrule
		Calendário e progresso & Os artefatos de cada entrega serão concluídos no prazo? & \textbf{I1} -- Comprometido \textit{versus} concluído: artefatos planejados para a entrega e artefatos integrados ao ramo principal (\%) & Faltando uma semana para a entrega, conclusão abaixo de 80\% exige redistribuição de tarefas \\
		& O trabalho restante está diminuindo no ritmo planejado? & \textbf{I2} -- \textit{Burndown}: \textit{issues} abertas do marco da entrega, por dia & Desvio maior que 10\% em relação à linha planejada é discutido na reunião seguinte \\
		Recursos e custos & O esforço está distribuído entre os membros? & \textbf{I3} -- Esforço por membro: horas registradas na planilha de esforço e participação percentual & Participação abaixo de 10\% ou acima de 40\% do total exige rebalanceamento \\
		Desempenho dos processos & A integração das fases é ágil? & \textbf{I4} -- Tempo de ciclo da fase: do primeiro \textit{commit} do ramo até a integração do PR (horas) & Ciclo maior que cinco dias indica integração tardia e risco de conflito \\
		& A revisão encontra problemas antes da integração? & \textbf{I5} -- Retrabalho: correções feitas após revisão sobre o total de alterações do PR (\%) & Retrabalho acima de 30\% motiva \textit{checklist} de revisão baseado nos critérios do roteiro \\
		\bottomrule
	\end{tabularx}
	\fonte{Elaborado pela equipe, com base em \citeonline{mcgarry2002} e \citeonline{psmcid2}.}
\end{table}

A primeira medição foi feita em 19/10/2026, a partir do histórico do repositório (\autoref{tab:psm-medicao}). Os dados brutos e o procedimento de extração estão na página de gestão do projeto, na GitPage.

\begin{table}[htb]
	\caption{Primeira medição dos indicadores de gestão (19/10/2026)}
	\label{tab:psm-medicao}
	\centering
	\scriptsize
	\begin{tabularx}{\textwidth}{L{2.6cm} L{4.2cm} Y}
		\toprule
		\textbf{Indicador} & \textbf{Valor medido} & \textbf{Interpretação e decisão} \\
		\midrule
		I1 -- Comprometido \textit{vs.} concluído (Entrega 1) & 14 de 17 artefatos planejados (82\%): 7 páginas da Fase 1 e 7 da Fase 2 integradas; artigo, tabela de contribuição e \textit{release} pendentes & Dentro do critério a sete dias da entrega. Decisão: priorizar o artigo e a tabela de contribuição, que dependem de todos os membros \\
		I2 -- \textit{Burndown} & Não mensurável: nenhuma \textit{issue} ou marco criado & Lacuna de processo. Decisão: criar os marcos E1 e E2 e uma \textit{issue} por artefato, para permitir o \textit{burndown} a partir da Fase 3 \\
		I3 -- Esforço por membro & Planilha ainda não adotada; como \textit{proxy}, \textit{commits} sem \textit{merge}: 9, 3, 3, 2 e 0 por membro & \textit{Commits} não medem esforço (reuniões e revisões não aparecem). Decisão: adotar a planilha de esforço, que também alimentará a tabela de contribuição \\
		I4 -- Tempo de ciclo & Fase 1: 49,0~h; Fase 2: 16,7~h & Ambos dentro do critério \\
		I5 -- Retrabalho & Não medido: os PRs não registraram rodadas de revisão & Decisão: registrar a revisão nos PRs a partir da Fase 3 \\
		\bottomrule
	\end{tabularx}
	\fonte{Elaborado pela equipe, a partir do histórico do repositório.}
\end{table}

\section{Aspectos éticos e rastreabilidade dos dados}
\label{sec:etica}

Nenhuma tarefa da avaliação altera dados acadêmicos: todas são consultas ou emissões de documentos feitas pelos próprios participantes, com as suas credenciais. A participação de estudantes em testes e no questionário é voluntária e precedida de termo de consentimento; as sessões de teste só são gravadas mediante consentimento, conforme a premissa 6 do roteiro \cite{ramos2026roteiro}. Os participantes são identificados apenas por código (P1, P2, \ldots) e por perfil, e capturas de tela têm os dados pessoais ocultados. Para garantir a rastreabilidade (premissa 5), os dados brutos anonimizados, os dados consolidados e os procedimentos de cálculo serão mantidos na pasta \texttt{dados/} do repositório; as gravações, por conterem dados pessoais, terão acesso restrito à equipe e à professora.

\section{Declaração de uso de inteligência artificial}
\label{sec:uso-ia}

A equipe usou o Claude Opus 5.5 (Claude Code) para: (a) analisar o roteiro e apoiar a redação e a formatação das páginas da Fase 1; (b) apoiar a redação da Fase 2 e verificar a consistência entre as páginas e os critérios do roteiro; e (c) estruturar este artigo no modelo LaTeX, condensar as páginas das Fases 1 e 2 e propor o questionário aplicado aos estudantes. O DeepSeek v4 foi usado na estruturação inicial do repositório MkDocs. Todo conteúdo gerado foi conferido pela equipe contra as páginas do repositório e o roteiro; as referências foram verificadas uma a uma, e os números citados (notas da matriz, faixas de pontuação e medições de gestão) foram conferidos com as fontes originais antes da integração. O registro completo, com datas e responsáveis, está na página de uso de IA da GitPage.
